\lecture{2}
\title{Nondeterminism, Closure Properties, Regular Expressions}

\begin{document}

\subsection{Nondeterministic Finite Automata}

We first introduce a new variant of finite automata:
the \textbf{nondeterministic finite automata} (or \textbf{NFA}).
\begin{center}
    \frame{\includegraphics[width=0.4\textwidth]{lecture-02/NFA-example}}
\end{center}
In spirit, it's a finite automaton that allows processes to fork;
if \emph{any} process ends at an accept state,
the automaton accepts.

\textbf{Definition.}
A nondeterministic finite automaton
is a finite automaton with the following alterations:
\begin{itemize}
    \item A state can have multiple outward edges
    with the same label.
    \item An outward edge can have the empty string $\epsilon$
    as a label.  (Traversing such an edge is free.)
\end{itemize}
Formally, the distinction is simply that
the transition function is now
$\delta \colon Q \times \Sigma_{\epsilon} \to \mathcal{P}(Q)$
(where $\Sigma_{\epsilon} := \Sigma \cup \{\epsilon\}$).

\begin{remark}
    There's an \textcolor{purple}{\emph{important}} note
    about both NFAs and DFAs that should be explained explicitly now.
    
    Sometimes, you might read in a character \str{c}
    while at a state with no outgoing edge with label \str{c}.
    In such a case, you go to a \textbf{dead state},
    meaning the string is automatically rejected.

    This is true even if the state you are at
    is the accept state.
    It is because of this that \str{aaaa} is a regular language,
    for example.
\end{remark}

\textbf{Example.}
In the example above:
\begin{itemize}
    \item If $M$ reads in \str{aba},
    then the end state could be any of $\{q_1, q_2, q_4, \textsc{Dead}\}$.
    Since $q_4$ is an accept state, $M$ accepts \str{aba}.
    \item If $M$ reads in \str{bbbbbaaaaaaa},
    then the end state could be any of $\{q_1, q_2, \textsc{Dead}\}$.
    So $M$ rejects \str{bbbbbaaaaaaa}.
    \item If $M$ reads in \str{aaabbb},
    then the end state could be any of $\{q_1, \textsc{Dead}\}$.
    So $M$ rejects \str{aaabbb}.
\end{itemize}

It turns out, very importantly,
that NFAs and DFAs are just different expressions of \emph{the same thing}.

\begin{theorem}{NFA Implies DFA}
    A language $A$ is recognized by some DFA
    if and only if $A$ is recognized by some NFA.
\end{theorem}

\begin{proof}
    If $A$ is recognized by some DFA $M$,
    then $A$ is recognized by the same FA $M$
    if interpreted as an NFA.

    Meanwhile, if $A$ is recognized by some NFA $N$,
    then we can design a DFA $M$ that recognizes $A$
    by taking $M$ to hold one state
    for each subset of the set of states of $N$.
    Details omitted; it's not very nontrivial. ~ $\blacksquare$
\end{proof}

So although it \emph{feels} that nondeterminism allows for
more generality than determinism\dots
it actually changes nothing!
The upshot, of course, is that in a lot of contexts,
NFAs are much easier to construct than DFAs.

\subsection{Regular Expressions are Regular Languages}

To demonstrate the strength of NFAs,
we'll now use them to prove the following theorem.

\begin{theorem}{Regularness}
    Every regular expression is a regular language.
\end{theorem}

It suffices to prove the closure of regular languages
under union, concatenation, and star.
NFAs do this pretty easily.

\begin{center}
    \frame{\includegraphics[width=0.3\textwidth]{lecture-02/union-proof}}
    \emph{(Union Proof)}
\end{center}

\begin{center}
    \frame{\includegraphics[width=0.5\textwidth]{lecture-02/concat-proof}}
    \emph{(Concatenation Proof)}
\end{center}

\begin{center}
    \frame{\includegraphics[width=0.3\textwidth]{lecture-02/star-proof}}
    \emph{(Star Proof)}
\end{center}

And that's all for today.

\end{document}